\section{Results}\label{sec:results}

\subsection{Detector configuration selection}\label{sec:detectorSelectionResults}

Based on the validation comparison in Table~\ref{tab:configurationSelection}, we selected Claude Opus 4.5 \nobreakcite{opus45Card}, with reasoning disabled, temperature \num{0.0}, and the Zero-Shot prompts \nobreakcite{DBLP:journals/csur/LiuYFJHN23}. The comparison also included GPT-5.2, GLM-5, and gpt-oss-120B \nobreakcite{gpt52,DBLP:journals/corr/abs-2602-15763,DBLP:journals/corr/abs-2508-10925}, with Few-Shot, Chain-of-Thought, and Tree-of-Thought-inspired prompts following established designs \nobreakcite{DBLP:conf/nips/BrownMRSKDNSSAA20,DBLP:conf/nips/KojimaGRMI22,DBLP:conf/nips/Wei0SBIXCLZ22,tree-of-thought-prompting,DBLP:conf/nips/YaoYZS00N23}. Three models reported the same displayed reference-support-weighted F1-score of \num{0.88} with Zero-Shot prompts. Claude Opus 4.5 completed that comparison in \num{367} seconds, the shortest observed runtime. The reasoning-oriented prompts reached \num{0.90} F1 with higher costs and runtimes. We therefore selected the simpler configuration for the held-out and corpus runs.

\begin{table}[tbp]
    \centering
    \scriptsize
    \caption{Originally reported validation results used to select the detector configuration. Precision, recall, and F1 are reference-support-weighted occurrence-agreement metrics over \num{20} validation repositories. Costs include value-added tax, and runtime is in seconds. The first four rows compare models with Zero-Shot prompts; the final three compare alternative prompts with Claude Opus 4.5.}
    \label{tab:configurationSelection}
    \begin{tabular}{llrrrrr}
        \hline
        \textbf{Model} & \textbf{Prompt} & \textbf{P} & \textbf{R} & \textbf{F1} & \textbf{Cost (USD)} & \textbf{Time (s)} \\
        \hline
        GPT-5.2 & Zero-Shot & 0.85 & 0.92 & 0.88 & 26.16 & 2423 \\
        GLM-5 & Zero-Shot & 0.88 & 0.89 & 0.88 & 11.77 & 5774 \\
        Claude Opus 4.5 & Zero-Shot & 0.88 & 0.89 & 0.88 & 29.97 & 367 \\
        gpt-oss-120B & Zero-Shot & 0.83 & 0.87 & 0.83 & 1.49 & 2808 \\
        Claude Opus 4.5 & Few-Shot & 0.88 & 0.89 & 0.88 & 45.00 & 352 \\
        Claude Opus 4.5 & Chain-of-Thought & 0.92 & 0.90 & 0.90 & 40.86 & 956 \\
        Claude Opus 4.5 & Tree-of-Thought & 0.92 & 0.90 & 0.90 & 81.43 & 4701 \\
        \hline
    \end{tabular}
\end{table}

\subsection{RQ1: What occurrence-level agreement does the selected detector achieve on held-out jOOQ projects?}\label{sec:rq1Agreement}

In one run, the selected detector returned \num{536} predictions from \num{502} relevant files in the \num{20}-repository test partition. The original reference contained \num{523} spans. At the prespecified intersection-over-union threshold of \num{0.50}, the one-to-one matching procedure yielded \num{460} true positives, \num{76} false positives, and \num{63} false negatives. The corresponding micro precision, recall, and F1-score were \num{0.858}, \num{0.880}, and \num{0.869}.

Table~\ref{tab:toolDetections} shows that the aggregate result masks substantial class differences. F1-scores ranged from \num{0.481} for Fear of the Unknown to \num{0.974} for 31 Flavors. Implicit Columns, 31 Flavors, and Rounding Errors each exceeded \num{0.89}, whereas Keyless Entry, Fear of the Unknown, and Poor Man's Search Engine remained below \num{0.70}. The macro precision, recall, and F1-score were \num{0.797}, \num{0.813}, and \num{0.793}; the reference-support-weighted values were \num{0.885}, \num{0.880}, and \num{0.877}. These aggregates reinforce the need to interpret each corpus flag class in light of its own held-out agreement.

\begin{table}[H]
    \centering
    \small
    \caption{Occurrence-level agreement with the original reference spans at intersection over union \num{0.50}. True positives (TP), false positives (FP), and false negatives (FN) are one-to-one event-matching counts; true negatives are undefined for this localisation task. The final row pools all classes.}
    \label{tab:toolDetections}
    \begin{tabular}{lrrrrrr}
        \hline
        \textbf{Antipattern} & \textbf{TP} & \textbf{FP} & \textbf{FN} & \textbf{Precision} & \textbf{Recall} & \textbf{F1} \\
        \hline
        Implicit Columns & 253 & 4 & 17 & 0.984 & 0.937 & 0.960 \\
        ID Required & 89 & 13 & 16 & 0.873 & 0.848 & 0.860 \\
        Keyless Entry & 35 & 30 & 1 & 0.538 & 0.972 & 0.693 \\
        Fear of the Unknown & 19 & 24 & 17 & 0.442 & 0.528 & 0.481 \\
        31 Flavors & 38 & 1 & 1 & 0.974 & 0.974 & 0.974 \\
        Poor Man's Search Engine & 13 & 4 & 8 & 0.765 & 0.619 & 0.684 \\
        Rounding Errors & 13 & 0 & 3 & 1.000 & 0.813 & 0.897 \\
        \textbf{Micro total} & \textbf{460} & \textbf{76} & \textbf{63} & \textbf{0.858} & \textbf{0.880} & \textbf{0.869} \\
        \hline
    \end{tabular}
\end{table}

\subsubsection{Class-specific errors}\label{sec:detectionErrors}

The disagreements in Table~\ref{tab:detectionErrors} cluster around local query semantics, inferred relationships, missing-value conventions, and occurrence grouping. Inspection also exposed omissions and span errors in the original reference annotations. Because detector disagreements guided this review, the primary result retains the original reference.

\begin{table}[H]
    \centering
    \small
    \caption{Dominant disagreement mechanisms observed during qualitative review of the held-out run. Counts refer to false positives or false negatives against the original reference spans.}
    \label{tab:detectionErrors}
    \begin{tabular}{p{2.9cm}p{7.8cm}}
        \hline
        \textbf{Antipattern} & \textbf{Observed disagreement mechanisms} \\
        \hline
        Implicit Columns & Fourteen of \num{17} false negatives used \texttt{returning()} or \texttt{returningResult()}. Other disagreements involved span boundaries and reference errors. \\
        ID Required & Six false negatives involved primary-key columns named \texttt{id}. One false positive proposed an unsuitable natural key; other disagreements exposed reference errors. \\
        Keyless Entry & Fifteen of \num{30} false positives involved audit-log tables, which the prompt did not exempt. One involved an ambiguous column name, and the remainder included missing or disputed reference annotations. \\
        Fear of the Unknown & All \num{17} false negatives involved defaults or sentinel values. Of \num{24} false positives, \num{23} exposed omissions in the original reference and one suggested a refactoring rather than an occurrence. \\
        31 Flavors & One false negative involved a numeric \texttt{CHECK} range encoding categories. The single false positive inferred a fixed value set from a \texttt{status} column without supporting schema evidence. \\
        Poor Man's Search Engine & Seven false negatives and two false positives resulted from grouping consecutive occurrences into one span. Two false positives used \texttt{LIKE} without wildcards, and one false negative placed \texttt{containsIgnoreCase} in an intermediate collection. \\
        Rounding Errors & All three false negatives involved \texttt{DOUBLE} coefficients later used to calculate monetary winnings. \\
        \hline
    \end{tabular}
\end{table}

Four mechanisms dominated: all-column \texttt{returning()} calls, keyless audit logs, sentinel defaults, and grouped search expressions. The remaining disagreements mainly involved line boundaries, sparse schema evidence, or class-specific semantics already summarised in the table.

The detector-informed revision contained \num{563} reference spans and changed the pooled totals to \num{512} true positives, \num{24} false positives, and \num{51} false negatives. These totals yield micro precision \num{0.955}, recall \num{0.909}, and F1-score \num{0.932}. We treat this result as an optimistic sensitivity analysis.

Micro F1 varied by \num{0.010} across the tested span boundaries. Table~\ref{tab:iouSensitivity} reports the complete threshold grid.

\begin{table}[tbp]
    \centering
    \small
    \caption{Sensitivity of pooled occurrence-level agreement to the line-span intersection-over-union (IoU) threshold. True positives (TP), false positives (FP), and false negatives (FN) follow the same one-to-one matching rule for all \num{536} predictions and \num{523} original reference spans. Micro F1 ranges from \num{0.863} to \num{0.873}.}
    \label{tab:iouSensitivity}
    \begin{tabular}{rrrrrrr}
        \hline
        \textbf{IoU} & \textbf{TP} & \textbf{FP} & \textbf{FN} & \textbf{Precision} & \textbf{Recall} & \textbf{F1} \\
        \hline
        0.25 & 462 & 74 & 61 & 0.862 & 0.883 & 0.873 \\
        0.50 & 460 & 76 & 63 & 0.858 & 0.880 & 0.869 \\
        0.75 & 457 & 79 & 66 & 0.853 & 0.874 & 0.863 \\
        1.00 & 457 & 79 & 66 & 0.853 & 0.874 & 0.863 \\
        \hline
    \end{tabular}
\end{table}

Only Implicit Columns and Poor Man's Search Engine changed across these thresholds. The pooled score remained within the \num{0.863}--\num{0.873} range.

A \num{10000}-replicate project-cluster bootstrap placed micro precision between \num{0.798} and \num{0.931}, micro recall between \num{0.820} and \num{0.913}, and micro F1 between \num{0.820} and \num{0.912}. These percentile ranges describe sensitivity to the composition of the selected test projects.

\subsection{RQ2: How frequently does the selected detector flag each of the seven SQL antipatterns in the 602-repository corpus, and how are those flags distributed across eligible files and repositories?}\label{sec:rq2CorpusFlags}

The reconstructed corpus frame contains \num{17988} relevant file occurrences in \num{602} repositories, and the detector produced \num{15931} positive flags. At least one retained class was flagged in \num{7653} files, or \qty{42.5}{\percent}, and \num{601} repositories, or \qty{99.8}{\percent}.

\begin{table}[H]
    \centering
    \scriptsize
    \setlength{\tabcolsep}{2pt}
    \caption{Selected-detector output in the identified \num{602}-repository corpus. G denotes the \num{7226} generated-schema files and A the \num{10762} application/query files; All contains \num{17988} relevant file occurrences. Repository coverage and file breadth give counts followed by percentages in parentheses. The two Fear of the Unknown repository counts overlap in two repositories. Repetition is flags per flagged file. Pooled and repository-equal densities are flags per 100 eligible files; the latter gives each eligible repository equal weight and includes zero-flag repositories. These are detector-output measures rather than prevalence estimates.}
    \label{tab:corpusFlags}
    \begin{tabular}{lrrrrr}
        \hline
        \makecell[l]{\textbf{Antipattern}\\\textbf{and role}} & \textbf{Flags} & \makecell{\textbf{Repository}\\\textbf{coverage}} & \makecell{\textbf{File}\\\textbf{breadth}} & \makecell{\textbf{Repeti-}\\\textbf{tion}} & \makecell{\textbf{Density}\\\textbf{pooled / repo.}} \\
        \hline
        Implicit Columns (A) & 7289 & 535 (88.9) & 2607 (24.2) & 2.80 & 67.73 / 92.63 \\
        ID Required (G) & 3597 & 523 (86.9) & 3591 (49.7) & 1.00 & 49.78 / 59.90 \\
        Keyless Entry (G) & 2153 & 200 (33.2) & 1269 (17.6) & 1.70 & 29.80 / 17.76 \\
        Fear of the Unknown (G) & 1240 & 197 (32.7) & 675 (9.3) & 1.84 & 17.16 / 20.59 \\
        Fear of the Unknown (A) & 16 & 10 (1.7) & 11 (0.1) & 1.45 & 0.15 / 0.17 \\
        31 Flavors (G) & 390 & 87 (14.5) & 292 (4.0) & 1.34 & 5.40 / 4.70 \\
        \makecell[l]{Poor Man's Search Engine (A)} & 583 & 114 (18.9) & 319 (3.0) & 1.83 & 5.42 / 4.27 \\
        Rounding Errors (G) & 663 & 86 (14.3) & 343 (4.7) & 1.93 & 9.18 / 7.63 \\
        \textbf{Any retained class (All)} & \textbf{15931} & \textbf{601 (99.8)} & \textbf{7653 (42.5)} & \textbf{2.08} & \textbf{88.56 / 93.63} \\
        \hline
    \end{tabular}
\end{table}

Table~\ref{tab:corpusFlags} separates class totals, repository coverage, flagged-file breadth, within-file repetition, and density in the corresponding source-role frame. The two most frequent classes have different output structures. ID Required flags cover \num{3591} of \num{7226} generated-schema files, or \qty{49.7}{\percent}, and average \num{1.00} per flagged file. Implicit Columns flags cover \num{2607} of \num{10762} application/query files, or \qty{24.2}{\percent}, and average \num{2.80} flags per flagged file. The first class is therefore broader across its eligible file frame, whereas the second repeats more within flagged files. Implicit Columns and ID Required together account for \num{10886} flags, or \qty{68.3}{\percent} of the output.

Across all repositories, pooled density is \num{88.56} flags per 100 relevant files. The repository-equal mean is \num{93.63}, and the median is \num{83.46}, with an interquartile interval of \numrange{51.28}{120.00}. Raw flag count correlates with relevant-file count at Spearman $\rho=\num{0.778}$. Large repositories contain \qty{43.8}{\percent} of relevant files and \qty{39.8}{\percent} of flags. Their broad density is \num{80.50}, below the small- and medium-repository densities of \num{93.19} and \num{96.65}.

Repository size therefore explains part of the raw concentration. The 61 highest-count repositories contain \qty{59.1}{\percent} of all flags and \qtyrange{48.1}{49.6}{\percent} of relevant files across the three-way cutoff tie. Their density is \numrange{1.47}{1.56} times that of the remaining repositories. High-count repositories combine greater file volume with higher density, while the large stratum as a whole has the lowest pooled density.

Exact-content sensitivity preserves both headline results. Of the \num{17988} file occurrences, \num{1109}, or \qty{6.17}{\percent}, belong to \num{488} exact-duplicate groups and contain \num{1080}, or \qty{6.78}{\percent}, of the raw flags. The groups contain \num{621} copies beyond their first occurrences. Taking the union of class-and-span events across byte-identical copies yields \num{15355} exact-content events over \num{17367} unique contents. Eight groups have inconsistent event sets, but selecting one physical representative from each group yields \numrange{15340}{15349} events without changing the interpretation. ID Required remains broader with almost no repetition, and Implicit Columns remains narrower and repetitive. The highest-count decile contains \qty{59.4}{\percent} of exact-content events, \qtyrange{50.1}{50.3}{\percent} of eligible unique-content weight, and \numrange{1.44}{1.46} times the remaining density. Byte-identical copies therefore do not explain the conclusions.

\subsection{RQ3: Which recurring source-fragment patterns are most frequent among detector flags for Implicit Columns and Poor Man's Search Engine in the 602-repository corpus?}\label{sec:rq3ApiPatterns}

Table~\ref{tab:apiManifestations} groups recurring source-fragment patterns among Implicit Columns and Poor Man's Search Engine flags. The two leading Implicit Columns patterns both selected all fields from a table. Together, they contained \num{5227} of \num{7289} flags, or \qty{71.7}{\percent}. Explicit field-array and asterisk patterns contributed a further \num{945} flags, or \qty{13.0}{\percent}. The reported recurring categories outside the combined catch-all accounted for \qty{94.1}{\percent} of the class flags.

For Poor Man's Search Engine, fragments matching \texttt{like(} contained \num{273} of \num{583} flags, or \qty{46.8}{\percent}. The next three patterns matched case-insensitive or contains-style method names. These four categories contained \num{517} flags, or \qty{88.7}{\percent}. The reported recurring categories outside the combined catch-all accounted for \qty{93.7}{\percent}.

\begin{table}[H]
    \centering
    \small
    \caption{Recurring post hoc source-fragment categories among detector flags for two query antipatterns. Categories are precedence-ordered string or regular-expression matches over detector-returned code fragments. Shares use the class flag total as denominator; they do not measure API-specific risk.}
    \label{tab:apiManifestations}
    \begin{tabular}{p{2.2cm}p{5.3cm}rr}
        \hline
        \textbf{Flag class} & \textbf{Source-fragment category} & \textbf{Flags} & \textbf{Share (\%)} \\
        \hline
        \makecell[l]{Implicit\\Columns} & \texttt{selectFrom} fragment & 3417 & 46.9 \\
        & \texttt{select()} fragment & 1810 & 24.8 \\
        & \texttt{fields()} fragment & 665 & 9.1 \\
        & \texttt{asterisk()} fragment & 280 & 3.8 \\
        & Other categorised patterns & 688 & 9.4 \\
        & Sparse matched or uncategorised & 429 & 5.9 \\
        \hline
        \makecell[l]{Poor Man's\\Search Engine} & \texttt{like(} fragment & 273 & 46.8 \\
        & \texttt{containsIgnoreCase(} fragment & 96 & 16.5 \\
        & \texttt{likeIgnoreCase(} fragment & 87 & 14.9 \\
        & \texttt{contains(} fragment & 61 & 10.5 \\
        & Other categorised patterns & 29 & 5.0 \\
        & Sparse matched or uncategorised & 37 & 6.3 \\
        \hline
    \end{tabular}
\end{table}
